\paragraph{再论\trBrane。}\label{sub:branes-revisited}
在 $(V,I_V)$ 上取定一个二次复\trVolume{}形式，意即复线之间的一个\trIsomorphism
\begin{equation*}
\eta_V^2\colon\lambda_{\mathbb{C}}^{\mathrm{top}}(V)^{\otimes 2}\longrightarrow\CC.
\end{equation*}
\begingroup\tolerance=9999\emergencystretch=3em\hbadness=10000
\trEquivalence{}地说，可以把它想成一个在相差正负号的意义下确定的复\trVolume{}形式\penalty0%
$\pm\eta_V\colon\lambda_{\mathbb{C}}^{\mathrm{top}}(V)\to\CC$。与之相伴的\textbf{\trSquaredPhase{}映射}%
（squared phase map，在特殊\trLagrange{}几何中人所共知）是\par\endgroup
\begin{equation}\label{eq:squared-phase-map}
\alpha_V\colon Gr(V)\longrightarrow S^1,\qquad
\alpha_V(\Lambda)=
\frac{\eta_V(v_1\wedge\cdots\wedge v_n)^2}
     {|\eta_V(v_1\wedge\cdots\wedge v_n)|^2}
\end{equation}
其中 $v_1,\ldots,v_n$ 是 $\Lambda$ 的任意一组\trBasis。这是\trWellDefined{}的，因为非零元素
$v_1\wedge\cdots\wedge v_n\in\lambda_{\mathbb{C}}^{\mathrm{top}}(V)$ 在相差乘以 $\RR^*$ 的意义下不依赖于%
\trBasis{}的选取，而平方便消去了正负号的歧义。我们进而把一个\textbf{\trLinear\trLagrange\trBrane}%
（linear Lagrangian brane）定义为三元组 $(\Lambda,\alpha^\#,P^\#)$，由一个\trLagrange{}子空间
$\Lambda$、一个满足 $\exp(2\pi i\,\alpha^\#)=\alpha_V(\Lambda)$ 的实数 $\alpha^\#$，以及一个主齐性%
$Pin_n$-空间 $P^\#$ 连同\trIsomorphism{} $P^\#\times_{Pin_n}\RR^n\cong\Lambda$ 组成。我们称
$\alpha^\#$ 为\textbf{\trRealValuedPhase}（real-valued phase）或\trGrading。
我们还需要考虑这类膜的\trFamily。也就是说，设 $E\to B$ 是一个满足 $2c_1(E)=0$
的\trSymplectic\trVectorBundle。这项条件保证：一旦取定一个\trCompatible{}的%
\trComplexStructure{} $I_E$，复\trLineBundle{} $\lambda_{\mathbb{C}}^{\mathrm{top}}(E)^{\otimes 2}$ 便是
平凡的，从而 $E$ 容许一个\trFibrewise{}的二次复\trVolume{}形式。取定这样一个形式
$\eta_E^2$，并令 $\alpha_E$ 是\trLagrangianGrassmannian{}丛 $Gr(E)\to B$ 上相伴的%
\trSquaredPhase{}映射。对一个\trLagrange\trSubbundle{} $F\subset E$ 而言，一个\trBraneStructure{}由
一个满足 $\exp(2\pi i\,\alpha^\#(b))=\alpha_E(F_b)$ 的\trSmooth{}函数
$\alpha^\#\colon B\to\RR$ 和 $F$ 上的一个 \trPinStructure{} $P^\#$ 组成。显然，存在%
\trBraneStructure{}的\trObstruction{}是：由\trSquaredPhase{} $b\mapsto\alpha_E(F_b)$
所表示的类 $\mu_F\in H^1(B)$，以及 $w_2(F)\in H^2(B;\ZZ/2)$。若这些\trObstruction{}消失，
则 $F$ 上诸\trBraneStructure{}的\trIsomorphism{}类构成
$H^0(B)\times H^1(B;\ZZ/2)$ 上的一个仿射空间。它由两种操作生成：用一个%
\trLocallyConstant{}的
$\ZZ$-值函数改变\trGrading，以及把 \trPinStructure{}与一个\trRealLineBundle{}作张量。回到
初始数据，我们应当指出，$E$ 的\trSubbundle{}的\trBraneStructure{}这一概念依赖于 $\eta_E^2$。更确切
地说，重要的只是它作为 $\lambda_{\mathbb{C}}^{\mathrm{top}}(E)^{\otimes 2}$ 的\trTrivialization{}的%
\trHomotopy{}类（若使 $\eta_E^2$ 发生\trDeformation，诸\trSquaredPhase{}函数会连续变化，
而可以把这一变化\trLift{}为诸\trRealValuedPhase{}的\trHomotopy），故在选取上有实效为
$H^1(B)$ 的自由度。
\begin{remark}\label{thm:grading-relative-orientation}
设 $\xi_E\subset\lambda_{\mathbb{C}}^{\mathrm{top}}(E)$ 是 $\pm\eta_E$ 取实值的一维实子空间。$\xi_E$
的非平凡性是寻找 $\eta_E^2$ 的整体平方根的\trObstruction。对任意 $b\in B$ 与 $F_b$
的任意一组\trBasis{} $\{v_k\}$，有 $v_1\wedge\cdots\wedge v_n\in\alpha_E(F_b)^{1/2}\xi_E$。假设%
\trBraneStructure{}已给定，便可把它乘以 $\exp(-\pi i\,\alpha^\#)$，得到一个\trIsomorphism{}
$\lambda^{\mathrm{top}}(F)\cong\xi_E$。这番讨论的要点是：\trGrading{}确定 $F$
\emph{相对于} $\xi_E$ 的一个\trOrientation。在\trGrading{}上添加 $\sigma\in\ZZ$，会
使这个\trOrientation{}乘以 $(-1)^\sigma$。
\end{remark}
在建立起这个一般框架之后，我们退回到一个更有限的情形，即平凡\trVectorBundle{}
$E=B\times V$ 连同常值的 $I_E=I_V$、$\eta_E^2=\eta_V^2$ 的情形。在这一语境下，眼前的
概念与先前定义的抽象膜之间是什么关系？容易看出，映射\eqref{eq:squared-phase-map}
表示\trMaslovClass{} $\mu_V$。此外，\trShortExactSequence{} $\ZZ/2\to Pin_n\to O_n$
诱导一个\trFibration{} $BPin_n\to BO_n$，其\trFiber{}为 $B\ZZ/2=\RP^\infty$，它可以等同
于在表示 $w_2$ 的映射 $BO_n\to K(\ZZ/2,2)$ 之下\trUniversal\trFibration{}的%
\trPullback（这只不过
是 $w_2$ 与 \trPinStructure{}之间关系的一个\trHomotopy{}论式的重新表述）。由此可见，$Gr(V)^\#$
可以取作\trPullback{} $\alpha_V^*(\RR\to S^1)$ 与 $i_V^*(BPin_n\to BO_n)$ 的%
\trFibreProduct，其中 $i_V$ 的定义是：先把 $Gr(V)$ \trEmbedding{}到由 $n$ 平面构成的
通常\trGrassmannian{}中，再\trEmbedding{}到 $BO_n$ 中。于是，对给定的 $F\subset B\times V$，
相应映射 $B\to Gr(V)$ 到 $Gr(V)^\#$ 的诸\trLift{}的\trHomotopy{}类，与诸二元组
$(\alpha^\#,P^\#)$ 的\trIsomorphism{}类之间有一一对应。这在\trHomotopy{}的层面上
回答了我们的问题；这在很大程度上令人满意，但并非完全如此。然而，与其寻求一个形式上
更完美的表述，我们觉得不如快速回顾一下主要应用来得更有启发性，因为那将具体地表明，
\trBraneStructure{}的新定义在功能上与旧定义\trEquivalence。
取一对\trLinear\trLagrange\trBrane{} $\Lambda_k^\#=(\Lambda_k,\alpha_k^\#,P_k^\#)$
（$k=0,1$），设 $(\Lambda_0,\Lambda_1)$ \trTransverse{}相交。正如\cref{sub:abstract-brane-structures}
中那样，可以关联于它们一个绝对\trIndex{}与一个\trOrientation{}空间。具体地说，取一条
路径 $\lambda\in\Omega^-(Gr(V),\Lambda_0,\Lambda_1)$ 以及关联丛
$F_\lambda\subset[0;1]\times V$ 上的一个\trBraneStructure{} $(\alpha_\lambda^\#,P_\lambda^\#)$，使%
\trGrading{}满足 $\alpha_\lambda^\#(k)=\alpha_k^\#$，连同\trIsomorphism{}
$f_k\colon(P_\lambda^\#)_k\cong P_k^\#$（$k=0$ 与 $k=1$ 两处皆然）。其余的恰如从前
\eqref{eq:absolute-index-orientation-space}。或许值得把一些细节想清楚，以厘清%
\trBraneStructure{}究竟在何处介入。设给了某条路径 $\lambda$。存在唯一的函数
$\alpha_\lambda^\#\colon[0;1]\to\RR$ 满足 $\alpha_\lambda^\#(0)=\alpha_0^\#$ 与
$\exp(2\pi i\,\alpha_\lambda^\#(s))=\alpha_V(\lambda(s))$。按照我们的规定，它还需满足
$\alpha_\lambda^\#(1)=\alpha_1^\#$，这是对 $\lambda$ 相对两端点的\trHomotopy{}类的一个限制
（事实上，这挑出一个取定的\trHomotopy{}类）。与之相对，在 $F_\lambda$ 上配备%
\trPinStructure{}没有任何\trObstruction{}；然而，当继而考虑路径的单参数\trFamily{}时，
诸关联\trVectorBundle{}在 $[0;1]^2$ 上相对边界存在 \trPinStructure{}这件事，确实会对%
\trHomotopy{}类施加一个非平凡的限制，这保证\trOrientation{}空间不会出现非平凡的%
\trMonodromy。
\begin{remark}\label{thm:brane-orientation-reversal}
实际上，重要的只是 $(P_\lambda^\#,f_0,f_1)$ 的\trIsomorphism{}类，这意味着实质上
只有两种可能。要从一个过渡到另一个，只需把 $f_0$ 或 $f_1$ 乘以非平凡元素
$-e\in\ker(Pin_n\to O_n)$，这会诱导\trOrientation{}空间的一个\trOrientationReversing{}
的\trAutomorphism。
\end{remark}
\begin{example}\label{thm:pin1-example}
让我们看看最低的维数 $n=1$。考虑 $\RR$，取标准\trGenerator{} $v=1$。按定义
$Pin_1=\{\pm e,\pm v\}\subset Cl(\RR)=\RR[v]/(v^2-1)$，故\trProjection{} $Pin_1\to O_1$
是\trSplitting{}的；我们一劳永逸地选定像为 $\{e,v\}$ 的那个\trSplitting。于是，实\trLineBundle{}
$F\to B$ 上的一个 $Pin_1$-结构等同于在底空间 $B$ 上再选取另一个\trRealLineBundle。
特别地，一个\trLinear\trLagrange\trBrane{}由一条实直线 $\Lambda\subset V$、一个%
\trRealValuedPhase{} $\alpha^\#$ 以及另一条抽象实直线 $\beta$ 组成。对两个这样%
\trTransverse{}相交的膜，
\begin{equation}\label{eq:pin1-index-orientation}
\begin{aligned}
i(\Lambda_0^\#,\Lambda_1^\#)&=[\alpha_1^\#-\alpha_0^\#]+1,\\
o(\Lambda_0^\#,\Lambda_1^\#)&\cong\Hom(\beta_0,\beta_1)
\otimes\Lambda_1^{\otimes\,i(\Lambda_0^\#,\Lambda_1^\#)}.
\end{aligned}
\end{equation}
这里 $[\,\cdot\,]$ 表示下一个较低的整数。这不过是选取合适的路径并把
\eqref{eq:absolute-index-orientation-space} 展开写出的结果。
\end{example}
